\chapter{RUANG BANACH}

\section{Transformasi Alami}
 Ingat bahwa jika $V$ suatu ruang linear bernorma, maka
 \index{dual!ruang}%
 \index{ruang!dual}%
 \index{<unopurstar@$V^*$ (ruang dual dari suatu ruang linear bernorma)}%
\emph{ruang dual} $V^*$ adalah himpunan semua fungsional linear kontinu pada~$V$.  Pada $V^*$, penjumlahan dan
perkalian skalar didefinisikan secara titik demi titik. Normanya adalah norma operator biasa: jika $f \in V^*$, maka
 \[ \norm{f} := \inf\{M > 0\colon \abs{f(x)} \le M\norm{x}\quad \text{untuk setiap $x \in V$}\} = \sup\{\abs{f(x)}\colon \norm{x} \le 1\}. \]
Di bawah metrik yang diinduksi oleh norma ini, ruang dual $V^*$ lengkap dan karena itu merupakan ruang Banach.

\begin{defn}\label{defn_nls_adjoint} Jika $T\colon V \sto W$ suatu pemetaan linear kontinu antara ruang-ruang linear bernorma,
 \index{<unopurstar@$T^*$ (adjoin ruang linear bernorma)}%
$T^*\colon W^* \sto V^*$ didefinisikan seperti pada ruang vektor: $T^*(g) = g\,T$ untuk setiap
$g \in W^*$.  Pemetaan $T^*$ disebut
 \index{adjoin!pemetaan linear!antara ruang linear bernorma}%
 \index{linear!pemetaan!adjoin dari}%
\df{adjoin} dari~$T$.
\end{defn}

\begin{exam}\label{exam_dual_functor} Adjoin $T^*$ dari suatu pemetaan linear terbatas antara ruang-ruang linear bernorma $V$ dan
$W$ sendiri merupakan pemetaan linear terbatas antara $W^*$ dan~$V^*$. Dalam kategori ruang linear bernorma dan pemetaan linear
 \index{funktor!dualitas}%
 \index{funktor dualitas}%
kontinu, pasangan pemetaan $V \mapsto V^*$ (pemetaan objek) dan $T \mapsto T^*$ (pemetaan morfisme) merupakan suatu funktor
kontravarian dari $\cat{NLS_\infty}$ ke dirinya sendiri.
\end{exam}

Perhatikan bahwa Definisi~\ref{defn_nls_adjoint} berbeda dari definisi adjoin suatu pemetaan
antara ruang-ruang Hilbert (lihat~\ref{def000926}).  Karena setiap ruang Hilbert merupakan ruang linear bernorma, hal ini dapat saja
menimbulkan kebingungan.  Ada baiknya kita meyakinkan diri tentang fakta sederhana berikut: karena anti-isomorfisme (dan akibatnya
isomorfisme) antara ruang Hilbert dan dualnya yang diberikan oleh \emph{teorema Riesz--Fr\'echet} (lihat~\ref{0022057}
dan~\ref{cor_Hsp_iso_dual}), adjoin ruang linear bernorma dari suatu pemetaan linear terbatas \emph{sangat} mirip (meskipun
tentu tidak identik) dengan adjoin ruang Hilbertnya.

\begin{exam}\label{C057717}
Dual $V^{**}$ dari dual $V^*$ suatu ruang linear bernorma $V$ adalah \df{dual kedua} dari~$V$
 \index{dual!kedua}%
dan $T^{**} := \bigl(T^*\bigr)^*$ merupakan pemetaan linear terbatas dari $V^{**}$ ke~$W^{**}$.  Pada kategori
$\cat{BAN_{\boldsymbol\infty}}$ yang terdiri atas ruang Banach dan transformasi linear terbatas,
 \index{kedua!dual!funktor untuk ruang linear bernorma}%
 \index{funktor!dual kedua!untuk ruang linear bernorma}%
\df{funktor dual kedua} mencakup pemetaan objek $B \mapsto B^{**}$ dan pemetaan morfisme $T \mapsto T^{**}$.
Funktor ini kovarian.
\end{exam}

\begin{defn}\label{defn_nat_transf} Misalkan $\cat A$ dan $\cat B$ kategori, serta $F$,$G \colon \cat A \sto \cat B$
funktor kovarian.  Suatu
 \index{alami!transformasi}%
 \index{transformasi!alami}%
\df{transformasi alami} dari $F$ ke $G$ adalah pemetaan $\tau$ yang memasangkan kepada setiap objek
$A$ dalam $\cat A$ suatu morfisme $\tau_A \colon F(A) \sto G(A)$ dalam $\cat B$ sedemikian sehingga
untuk setiap morfisme $\alpha \colon A \sto A'$ dalam $\cat A$, diagram berikut komutatif.
   \[ \CD
      F(A) @>  F(\alpha)>>  F(A')  \\
         @V\tau_AVV                @VV\tau_{A'}V  \\
      G(A) @>> G(\alpha)>  G(A')
   \endCD \]
Kita menyatakan transformasi semacam itu dengan $\tau \colon F \sto G$. (Definisi transformasi alami
antara dua funktor kontravarian seharusnya jelas: cukup balikkan
panah-panah horizontal dalam diagram sebelumnya.)

Suatu transformasi alami $\tau \colon  F \sto G$ disebut
 \index{alami!ekuivalensi}%
 \index{ekuivalensi!alami}%
\df{ekuivalensi alami} jika setiap morfisme $\tau_A$ merupakan isomorfisme dalam~$\cat B$.
\end{defn}

\begin{defn}\label{defn_nat_embed} Jika $V$ suatu ruang linear bernorma, maka pemetaan $\sbsb{\tau}V \colon V \sto V^{**} \colon x
\mapsto x^{**}$, dengan $x^{**}(f) = f(x)$ untuk setiap $f \in V^*$, disebut
 \index{tauv@$\sbsb{\tau}V$ (pembenaman alami ruang linear bernorma)}%
 \index{alami!pembenaman!untuk ruang linear bernorma}%
 \index{pembenaman!alami}%
\df{pembenaman alami} dari $V$ ke~$V^{**}$. (Tidak sepenuhnya jelas bahwa pemetaan ini
injektif. Kita akan menunjukkannya dalam Proposisi~\ref{C067137}.)
\end{defn}

\begin{exam}\label{exam_2dual_nat} Dalam kategori $\cat{BAN_{\boldsymbol\infty}}$ yang terdiri atas ruang Banach
dan transformasi linear terbatas, misalkan $I$ adalah funktor identitas dan $(\,\cdot\,)^{**}$ adalah
funktor dual kedua.  Maka
 \index{alami!pembenaman!sebagai funktor}%
 \index{funktor!pembenaman alami sebagai}%
 \index{taub@$\tau_B$ (pembenaman alami $B$ ke $B^{**}$)}%
\emph{pembenaman alami} $\tau$ merupakan transformasi alami dari $I$ ke~$(\,\cdot\,)^{**}$.
\end{exam}

\begin{exam} Jelaskan secara tepat apa yang dimaksud ketika orang mengatakan bahwa isomorfisme
    \[ \fml C(X \uplus Y) \cong \fml C(X) \times \fml C(Y) \]
merupakan ekuivalensi alami.  Kategori-kategori yang terlibat adalah $\cat{CpH}$ (ruang Hausdorff kompak
dan pemetaan kontinu) serta $\cat{BAN_\infty}$ (ruang Banach dan transformasi linear
terbatas).
\end{exam}

\begin{exam}\label{C057724} Pada kategori $\cat{CpH}$ yang terdiri atas ruang Hausdorff kompak dan
pemetaan kontinu, misalkan $I$ funktor identitas dan $\fml C^*$ funktor yang
memetakan ruang Hausdorff kompak $X$ ke ruang Banach $\fml (C(X))^*$, serta memetakan suatu
fungsi kontinu $\phi\colon X \sto Y$ antara ruang-ruang Hausdorff kompak ke
transformasi linear kontraktif $\fml C^*\phi = (\fml C\phi)^*$. Definisikan
 \index{evaluasi!pemetaan}%
\emph{pemetaan evaluasi} $E$ dengan
  \[ E_X\colon X \sto \fml C^*(X)\colon x \mapsto E_X(x) \]
dengan $(E_X(x))(f) = f(x)$ untuk setiap $f \in \fml C(X)$.  (Notasi ini dapat merepotkan.  Dalam praktik, biasanya orang menulis
$E_x$ saja untuk $E_X(x)$.  Jadi kita memperoleh $E_x(f) = f(x)$, yang tentu tampak lebih baik daripada $(E_X(x))(f) = f(x)$.  Dalam kebanyakan
konteks, penyederhanaan notasi ini tampaknya tidak mungkin menimbulkan kebingungan.)
 \begin{enumerate}
  \item[(a)]  Perhatikan bahwa definisi ini memang bermakna. Jika $x \in X$, maka $E_X(x)$ benar-benar
merupakan anggota~$\fml C^*(X)$.
  \item[(b)] Tunjukkan bahwa pemetaan $E$ membuat diagram berikut komutatif:
     \[ \xymatrix@+30pt{{X}\ar[r]^*+{\phi} \ar[d]_*+{E_X}
            & {Y} \ar[d]^*+{E_Y} \\
          {\fml C^*(X)} \ar[r]_*+{\fml C^*(\phi)} & {\fml C^*(Y)}   }
     \]
 \end{enumerate}
\end{exam}

Contoh sebelumnya mungkin terasa agak mengecewakan.  Tampaknya kita sedang
berusaha membentuk suatu transformasi alami antara funktor identitas (pada
$\cat{CpH}$) dan funktor $\fml C^*$.  Masalahnya, tentu saja, ialah kedua
funktor tersebut memiliki kategori sasaran yang berbeda.  Pertanyaan ini akan kita bahas kelak.
Pemetaan evaluasi $E_X\colon X \sto \fml C^*(X)$ tentu tidak surjektif.
Perhatikan, misalnya, bahwa setiap fungsional berbentuk $E_X(x)$
mempertahankan perkalian (selain penjumlahan dan perkalian skalar). Salah satu akibatnya
ialah bahwa semua fungsional semacam itu terletak pada sfer satuan $\fml C^*(X)$.
Dengan suatu topologi baru (yang disebut \emph{topologi lemah-bintang}---lihat~\ref{C067434}) pada
$\fml C^*(X)$, jangkauan $E_X$ ternyata merupakan ruang Hausdorff kompak yang secara alami
ekuivalen dengan $X$ sendiri.  Yang lebih luar biasa lagi, pada akhirnya kita akan menunjukkan bahwa jangkauan $E_X$
dapat dipandang sebagai ruang ideal maksimal dari aljabar~$\fml C(X)$. (Lihat~\futurexref{11.2.20}{000731}.)

\begin{prop}\label{C067137} Jika $V$ suatu ruang linear bernorma, maka pembenaman alami
$\sbsb{\tau}V \colon V \sto V^{**}\colon x \mapsto x^{**}$ (yang didefinisikan dalam~\ref{defn_nat_embed})
merupakan pemetaan linear isometrik.
\end{prop}

\begin{defn} Suatu ruang Banach $B$ disebut
 \index{refleksif!ruang Banach}%
\df{refleksif} jika pembenaman alami $\sbsb{\tau}B\colon B \sto B^{**}$ (yang didefinisikan dalam proposisi sebelumnya) surjektif.
\end{defn}

Contoh~\ref{exam_2dual_nat} menunjukkan bahwa dalam kategori ruang Banach dan pemetaan linear terbatas,
pemetaan $\tau\colon B \mapsto \sbsb{\tau}B$ merupakan transformasi alami dari
funktor identitas ke funktor dual kedua.  Dengan menggunakan \emph{teorema Hahn--Banach}, kita dapat
mengatakan lebih banyak.

\begin{prop} Dalam kategori ruang Banach refleksif dan pemetaan linear terbatas, pemetaan
$\tau\colon B \mapsto \sbsb{\tau}B$ merupakan ekuivalensi alami antara funktor identitas dan
funktor dual kedua.
\end{prop}

\begin{exam} Setiap ruang Hilbert bersifat refleksif.
\end{exam}

\begin{exer} Misalkan $V$ dan $W$ ruang-ruang linear bernorma, $U$ subhimpunan konveks takkosong dari $V$,
dan $f \colon U \sto W$.  Tanpa menggunakan bentuk apa pun dari \emph{teorema nilai rata-rata}, tunjukkan
bahwa jika $df = \vc 0$ pada $U$, maka $f$ konstan pada~$U$.  (Untuk definisi $df$, yaitu
\emph{diferensial} dari~$f$, lihat Bagian 13.1--2 dalam catatan saya~\cite{Erdman:2007}.)
\end{exer}

\begin{exam}\label{C067181} Misalkan $l_\infty$ ruang Banach dari semua barisan bilangan
real yang terbatas, dan $c$ subruang $l_\infty$ yang terdiri atas semua barisan $x = (x_n)$
sedemikian sehingga $\lim_{n \sto \infty} x_n$ ada. Definisikan
  \[ f\colon c \sto \R\colon x \mapsto \lim_{n\sto \infty} x_n. \]
Maka $f$ merupakan fungsional linear kontinu pada $c$ dan $\norm f = 1$.  Berdasarkan
\emph{teorema Hahn--Banach}, $f$ memiliki perluasan ke $\wh f \in {l_\infty}^*$ dengan
$\norm{\wh f} = 1$. Untuk setiap
subhimpunan $A$ dari $\N$, definisikan barisan $\bigl(x^A_n\bigr)$ dengan $x^A_n = \left\{%
\begin{array}{ll}
    1, & \hbox{jika $n \in A$;} \\
    0, & \hbox{jika $n \notin A$.} \\
\end{array}%
\right.$ Selanjutnya, definisikan $ \mu\colon \sfml P(\N) \sto \R \colon A \mapsto \wh
f\bigl(x^A_n\bigr)$. Maka $\mu$ merupakan fungsi himpunan aditif hingga, tetapi tidak aditif
terhitung.
\end{exam}













\section{Teorema Alaoglu}
\begin{notn} Misalkan $V$ suatu ruang linear bernorma, $M \subseteq V$, dan $F \subseteq V^\ast$.
Kita definisikan
 \begin{align*}
         M^\perp &:= \{f \in V^*\colon f(x)=0 \text{ untuk setiap $x \in M$}\} \\
         F_\perp &:= \{x \in V\colon f(x)=0 \text{ untuk setiap $f \in F$}\}
  \end{align*}
$M^\perp$ dapat dibaca ``M perp'' atau ``M perp atas'', sedangkan $F_\perp$ dibaca ``F perp'' atau ``F perp bawah''.
Himpunan $M^\perp$ adalah
 \index{anihilator}%
 \index{<unopur@$M^\perp$ (anihilator suatu himpunan)}%
\df{anihilator} dari~$M$, dan $F_\perp$ adalah
 \index{praanihilator}%
 \index{<unopvr@$F_\perp$ (praanihilator suatu himpunan)}%
\df{praanihilator} dari~$F$.
\end{notn}

\begin{prop}[Anihilator]  Misalkan $B$ suatu ruang Banach, $M \subseteq B$, dan $F \subseteq B^*$. Maka
 \begin{enumerate}[label=\rm{(\alph*)}]
  \item Jika $M \subseteq N \subseteq B$, maka $N^\perp \subseteq M^\perp$.
  \item Jika $F \subseteq G \subseteq B^*$, maka $G_\perp \subseteq F_\perp$.
  \item $M^\perp$ merupakan subruang linear tertutup dari $B^*$.
  \item $F_\perp$ merupakan subruang linear tertutup dari $B$.
  \item $M \subseteq M^\perp{}_\perp$.
  \item $F \subseteq F_\perp{}^\perp$.
  \item $M^\perp{}_\perp = \bigvee M$.
  \item $M^\perp = B^*$ jika dan hanya jika $M \subseteq \{0\}$.
  \item $F_\perp = B$ jika dan hanya jika $F \subseteq \{0\}$.
  \item Jika $M$ suatu subruang linear dari $B$, maka $M^\perp = \{0\}$ jika dan hanya jika $M$
padat dalam $B$.
  \item\label{Fpp_reflexive_case} Jika $B$ refleksif, maka $F_\perp{}^\perp = \bigvee F$.
 \end{enumerate}
\end{prop}

\begin{proof}[\emph{Petunjuk pembuktian}] Untuk (k), tunjukkan bahwa $\tau(F_\perp) = F^\perp$ (dengan $\tau$
pembenaman alami) dan karena itu $F_\perp{}^\perp = F^\perp{}_\perp$.  \ns
\end{proof}

\begin{exam} Misalkan $B$ suatu ruang Banach yang \emph{tidak} refleksif.  Misalkan $F = \ran \tau$ (dengan $\tau$
pembenaman alami).  Maka $F_\perp{}^\perp = B^{**} \ne \bigvee F$, sehingga jelas bahwa~\ref{Fpp_reflexive_case}
dalam proposisi sebelumnya tidak harus berlaku dalam kasus nonrefleksif.
\end{exam}

Hasil berikut adalah analog ruang Banach dari \emph{teorema dasar aljabar linear}~\ref{00016025} dan
Proposisi~\ref{prop_ker_adj_perp_ran_op}.

\begin{prop}\label{prop_FTLA_Bsp} Misalkan $T \in \ofml B(B,C)$, dengan $B$ dan $C$ ruang-ruang Banach. Maka
 \begin{enumerate}[label=\rm{(\alph*)}]
  \item $\ker T^\ast = (\ran T)^\perp$.
  \item $\ker T = (\ran T^*)_\perp$.
\vskip 3 pt
  \item $\clo{\ran T} = (\ker T^*)_\perp$.
\vskip 3 pt
  \item $\clo{\ran T^*} \subseteq (\ker T)^\perp$.
 \end{enumerate}
\end{prop}

\begin{exam} Suatu contoh operator ruang Banach $T$ yang membuat $\clo{\ran T^*}$ dan $(\ker T)^\perp$ berbeda
diberikan dalam Latihan 7, halaman 243 dari~\cite{BrownP:1970}.
\end{exam}

Contoh berikut tentu merupakan akibat langsung yang sepele dari Proposisi~\ref{00015032}, tetapi cobalah memberikan
pembuktian langsung yang sederhana.

\begin{exam}\label{X_l2ball_not_cpt} Bola satuan tertutup dalam $l_2$ tidak kompak.
\end{exam}

Topologi norma biasa pada ruang linear bernorma ternyata terlalu besar untuk beberapa penerapan.  Seperti telah kita lihat,
topologi ini memiliki begitu banyak himpunan terbuka sehingga bola satuan tertutup tidak kompak (kecuali dalam kasus berdimensi hingga). Di
sisi lain, terlalu sedikit himpunan terbuka (sebagai contoh ekstrem, topologi indiskret) akan merusak teori dualitas karena tidak
menyediakan cukup banyak fungsional linear kontinu.  \emph{Topologi lemah-bintang} (yang didefinisikan
di bawah) pada dual suatu ruang linear bernorma ternyata tepat ukurannya.  Topologi ini memiliki cukup banyak himpunan terbuka untuk menjamin
teori dualitas yang kuat dan cukup sedikit himpunan terbuka untuk membuat bola satuan tertutup kompak
(lihat \emph{teorema Alaoglu}~\ref{C067441}).

\begin{defn}\label{C067434} Misalkan $V$ suatu ruang linear bernorma. 
 \index{lemah!topologi!pada ruang bernorma}%
 \index{topologi!lemah!pada ruang linear bernorma}%
\df{topologi lemah} pada $V$ adalah topologi lemah yang diinduksi oleh unsur-unsur~$V^*$. (Jadi,
topologi lemah pada ruang dual $V^*$ adalah topologi lemah yang diinduksi oleh
unsur-unsur~$V^{**}$.) Ketika suatu jaring $(x_\lambda)$ konvergen secara lemah ke vektor $a$ dalam $V$,
kita menulis
 \index{<arrowright@$x_\lambda \to^w a$ (konvergensi lemah)}%
$x_\lambda \to^w a$. 
 \index{w*@$w^*$-topologi (topologi lemah-bintang)}%
 \index{topologi!lemah-bintang}%
\df{topologi-$w^*$} (dibaca \emph{topologi lemah-bintang}) pada ruang dual $V^*$ adalah
topologi lemah yang diinduksi oleh unsur-unsur $\ran \tau$, dengan $\tau$ pembenaman alami
dari $V$ ke~$V^{**}$.  Ketika suatu jaring $(f_\lambda)$ konvergen secara lemah-bintang ke vektor $g$
dalam $V^*$, kita menulis
 \index{<arrowright@$f_\lambda \to^{w^*} g$ (konvergensi lemah-bintang)}%
$f_\lambda \to^{w^*} g$. Perhatikan bahwa ketika suatu ruang Banach $B$ refleksif (khususnya, untuk ruang
Hilbert), topologi lemah dan lemah-bintang pada $B^*$ identik.
\end{defn}

\begin{prop} Misalkan $f \in V^*$, dengan $V$ suatu ruang Banach.  Untuk setiap $\epsilon > 0$ dan setiap
subhimpunan hingga $F \subseteq V$, definisikan
  \[ U(f;F;\epsilon)
           = \{g \in V^*\colon \abs{f(x) - g(x)} < \epsilon \text{ untuk setiap $x \in F$}\}. \]
Keluarga semua himpunan semacam itu merupakan basis bagi topologi-$w^*$ pada~$V^*$.
\end{prop}

Hasil berikut, \emph{teorema Alaoglu}, menyatakan kekompakan-$w^*$ dari bola satuan tertutup dalam
dual suatu ruang linear bernorma $V$ dan merupakan salah satu teorema yang paling berguna dalam analisis fungsional.  Karena itu,
pada awalnya mungkin agak mengejutkan bahwa pembuktian lengkap yang dijelaskan secara cermat
begitu sulit ditemukan.  Pembuktian standar pada dasarnya sangat sederhana; intinya hanyalah mengenali
bahwa anggota-anggota suatu keluarga $\fml F$ dari fungsi bernilai skalar dapat dipandang
  \begin{enumerate}
     \item[(i)] sebagai fungsi yang didefinisikan pada hasil kali takterhitung dari subhimpunan kompak medan skalar; atau
     \item[(ii)] sebagai anggota bola satuan tertutup dari~$V^*$.
  \end{enumerate}
Inti pembuktiannya adalah meyakinkan diri bahwa konvergensi suatu jaring fungsi semacam itu dalam ruang hasil kali
tepat ``sama'' dengan konvergensi-$w^*$ jaring tersebut dalam bola satuan tertutup dari~$V^*$.  Upaya untuk sekaligus
mengidentifikasi dan membedakan dua himpunan fungsi ini menantang secara notasi.  Jadi, saran terbaik saya ialah mengerjakannya
sendiri sampai Anda melihat betapa cerdik, tetapi pada dasarnya sederhana, argumen tersebut.

\begin{thm}[Teorema Alaoglu]\label{C067441} Bola satuan tertutup dalam dual suatu ruang linear bernorma
 \index{teorema Alaoglu}%
kompak dalam topologi-$w^*$.
\end{thm}

\begin{proof}[\emph{Petunjuk pembuktian}] Misalkan $V$ suatu ruang linear bernorma, dan nyatakan dengan $[V^*]_1$ bola satuan tertutup dari
ruang dualnya.  Untuk setiap $x \in V$, definisikan $D_x = \{\lambda \in \K \colon \abs\lambda \le \norm x\}$.  Hasil kali
\smash[b]{$\mathbf P = \prod\limits_{x \in V} D_x$} kompak berdasarkan \emph{teorema Tychonoff}.  Tinjau fungsi
   \[ \Phi\colon [V^*]_1 \sto \mathbf P \colon f \mapsto \bigl(f(x)\bigr)_{x \in V} \]
dengan $[V^*]_1$ dilengkapi topologi-$w^*$ dan $\mathbf P$ dilengkapi topologi hasil kali (lihat~\ref{C021437}).  Di sini kita menulis
$\bigl(f(x)\bigr)_{x \in V}$ untuk keluarga bilangan terindeks (atau barisan tergeneralisasi), bukan notasi yang lebih lazim
$\bigl(f_x\bigr)_{x \in V}$.  (Untuk pembahasan lebih lanjut tentang keluarga terindeks, lihat~\cite{Erdman:2007}, Bagian~7.2.)  Mudah dilihat bahwa
$\Phi$ injektif.  Untuk menunjukkan bahwa pemetaan itu kontinu, misalkan $\bigl(f_\lambda\bigr)_{\lambda \in \Lambda}$ suatu jaring dalam $[V^*]_1$
dan $g$ suatu fungsi dalam $[V^*]_1$ sedemikian sehingga $f_\lambda \to^{w^*} g$.  Perhatikan bahwa hal ini terjadi jika dan hanya jika
$\pi_x(f_\lambda) \sto \pi_x(g)$ untuk setiap $x \in V$ (dengan $\pi_x$ proyeksi koordinat pada~$\mathbf P$---lihat~\ref{C015531}).
Terakhir, misalkan $\bigl(f_\lambda\bigr)_{\lambda \in \Lambda}$ suatu jaring dalam $[V^*]_1$ sedemikian sehingga $\bigl(\Phi(f_\lambda)\bigr)_{\lambda \in \Lambda}$
konvergen dalam~$\mathbf P$.  Tunjukkan bahwa jaring $\bigl(f_\lambda(x)\bigr)_{\lambda \in \Lambda}$ konvergen dalam $\K$ untuk setiap $x \in V$.
Misalkan $g(x) = \lim\limits_\lambda f_\lambda(x)$ untuk setiap $x \in V$, lalu buktikan bahwa $g$ linear, $\norm g \le 1$,
$f_\lambda \to^{w^*} g$ dalam $[V^*]_1$, dan $\Phi(f_\lambda) \sto \Phi(g)$ dalam~$\mathbf P$.  \ns
\end{proof}

Sepengetahuan saya, versi pembuktian standar ini yang paling ramah pembaca terdapat dalam~\cite{BrownP:1977}, Teorema 15.11.
Pembuktian yang jauh lebih singkat, jauh lebih mudah secara notasi, jauh lebih elegan, dan jauh kurang mencerahkan tersedia dalam~\cite{Pedersen:1995}, Teorema 2.5.2.
Menurut saya, pembuktian itu kurang mencerahkan karena bergantung pada jaring universal.  Suatu jaring dalam himpunan $S$ disebut
 \index{universal!jaring}%
 \index{jaring!universal}%
\df{universal} jika untuk setiap subhimpunan $T$ dari $S$, jaring tersebut pada akhirnya berada dalam $T$ atau pada akhirnya berada dalam komplemennya. Meskipun
setiap jaring dapat dibuktikan memiliki subjaring universal, bahkan mencari satu contoh nontrivial dari jaring universal (yakni, jaring yang tidak
pada akhirnya konstan) memerlukan \emph{aksioma pilihan}.






